

%--------------------------------------------------------------------------------------
%--------------------------------------------------------------------------------------
%--------------------------------------------------------------------------------------
\section{Antecedentes}
\label{sec:antecedentes}
    La creciente necesidad de generar estadísticas confiables y detalladas a niveles de desagregación geográfica o demográfica menores a los inicialmente previstos por los diseños muestrales ha impulsado el desarrollo y la aplicación de los modelos de Estimación en Áreas Pequeñas, consolidándolos como una herramienta metodológica indispensable para la formulación de políticas públicas basadas en evidencia. El problema fundamental que resuelven los modelos SAE es la insuficiencia del tamaño muestral en los dominios de interés, que provoca que los estimadores directos carezcan de la precisión adecuada, presentando errores estándar inaceptablemente altos. La solución radica en el uso de modelos estadísticos que ``toman prestada fuerza'' de otras áreas y utilizan información auxiliar externa altamente correlacionada con el parámetro de interés, logrando predicciones más confiables. A lo largo del tiempo, la aplicación de estos modelos ha sido extensa y diversa, cubriendo desde indicadores económicos hasta variables latentes en la educación.
    %-XXXXXXXXXXXXXXXXXXXXXXXXXXXXXXXXXXXXXXXXXXXXXXXXXXXXXXXXXXXX
    % Esto se repite mucho: Estimación en Áreas Pequeñas (SAE, por sus siglas en inglés)
    % Ya se definió, por tanto solo se usa SAE.
    % Deben revisar el uso de las siglas, porque algunas o no se definen y otras se repiten como en este caso.
    %-XXXXXXXXXXXXXXXXXXXXXXXXXXXXXXXXXXXXXXXXXXXXXXXXXXXXXXXXXXXX

    Los antecedentes históricos comienzan con el fundamental Modelo Fay--Herriot (1979), desarrollado por Fay y Herriot para estimar el ingreso per cápita en áreas pequeñas de Estados Unidos \cite{BernalGonzalez2017}. Este modelo a nivel de área lineal mixto fue pionero al combinar el estimador directo con una regresión sintética, resolviendo el problema de la falta de fiabilidad en estimaciones directas para poblaciones con muestras limitadas.

    Posteriormente, la necesidad de medir la pobreza a nivel desagregado condujo al desarrollo de extensiones del Fay--Herriot y modelos alternativos. Como se menciona en \cite{Molina2019}, la CEPAL aplicó un Modelo Lineal Mixto (similar a Fay--Herriot) para estimar el los ingresos medios y las incidencias de probreza no extrema para las secciones censales y para los dos géneros en Montevideo, superando la escasez de información en municipios no cubiertos  por las fuentes tradicionales. Abordando problemas laborales, en \cite{MendozaZea2019} emplearon un Modelo Fay--Herriot Espacial para estimar la tasa de desempleo y el ingreso medio en Cundinamarca, Colombia. El uso de esta metodología, que incorpora efectos espaciales y variables auxiliares como puntajes de pruebas académicas, demostró ser más efectiva que las estimaciones directas, especialmente para el ingreso medio.

    Para problemas que involucran la estimación de proporciones (como tasas de pobreza, cuyo dominio está restringido al intervalo (0, 1)), se han desarrollado metodologías especializadas. Como se menciona en \cite{BernalGonzalez2017} se implementó un Estimador Completamente Bayesiano Beta (ECBB) con un enfoque Jerárquico de Bayes (HB). Esta metodología novedosa fue aplicada para estimar el porcentaje de hogares con carencias para la población infantil y juvenil en las UPZ de Bogotá, abordando el problema de variables cuya distribución está acotada entre 0 y 1. De manera similar, en Chile, en \cite{MarinTralmaGuerrero2024} aplicaron un Modelo EBLUP a Nivel de Área para estimar indicadores de victimización a nivel comunal, resolviendo el problema de las proporciones mediante una transformación arcoseno para asegurar que las estimaciones se mantuvieran dentro del rango.

    En el campo del rendimiento académico, que a menudo se vincula con la pobreza, las metodologías EAP se han combinado con técnicas de encuestas complejas para resolver el problema de datos faltantes y latentes. Como se menciona en \cite{TellezTrujilloPedraza2020} se enfocaron en la estimación de los coeficientes de regresión ('s) de un modelo lineal en áreas pequeñas (departamentos de Colombia) utilizando la técnica de Valores Plausibles (Imputación Múltiple) y un modelo Fay-Herriot Multivariado. Esta combinación resuelve el problema de las matrices de datos incompletas resultantes de diseños muestrales complejos, como el muestreo de bloques incompletos. Finalmente, en \cite{ColeyTellez2024} propusieron un modelo híbrido que integra la Teoría de Respuesta al Ítem (TRI) con Valores Plausibles y un Modelo Fay-Herriot Espacial. Este modelo se aplicó para estimar el rendimiento promedio de lectura en PISA 2018, utilizando variables auxiliares socioeconómicas como la tasa de incidencia de pobreza. Los resultados de este modelo demostraron una reducción en la variabilidad de las estimaciones en comparación con los datos oficiales.

    En Chile, la Estimación en Áreas Pequeñas ha sido adoptada oficialmente para la formulación de políticas públicas. El Ministerio de Desarrollo Social publicó una metodología para la estimación de la tasa de pobreza por ingresos a nivel comunal. Para resolver la imprecisión de las estimaciones directas de la Encuesta Casen, se empleó un Modelo Compuesto (tipo Fay--Herriot). El proceso implicó el suavizamiento de los factores de expansión por medio del método Potter, la estimación de la varianza sintética y la varianza directa, y la combinación de ambas a través del ponderador o $shrinkage\ factor$ \cite{INEChile2016}.

    %-XXXXXXXXXXXXXXXXXXXXXXXXXXXXXXXXXXXXXXXXXXXXXXXXXXXXXXXXXXXX
    % ADeben revisar la cita [9]
    %-XXXXXXXXXXXXXXXXXXXXXXXXXXXXXXXXXXXXXXXXXXXXXXXXXXXXXXXXXXXX

    Estos estudios demuestran que, ya sea para atacar el problema de la pobreza mediante modelos de micro-nivel (ELL) en Brasil, o para refinar las tasas de pobreza y seguridad utilizando extensiones del modelo Fay--Herriot (con $shrinkage\ factor$ y transformación arcoseno) en Chile, la Estimación en Áreas Pequeñas es una herramienta flexible y esencial para generar estadísticas fiables y desagregadas que orientan la política pública en estos países.
